\documentclass[a4paper,landscape,headrule,footrule]{foils}

%%
%%% macros for Theories of Grammar
%%%
\usepackage{polyglossia}
\setdefaultlanguage{english}
%\setmainfont{TeX Gyre Pagella}


\newcommand{\logo}{~}
\newcommand{\header}[3]{%
\title{\vspace*{-2ex} \large JMORF/MORS --- Morpho-Syntax
\\[2ex] \Large  \emp{#2} \\ \emp{#3}}
\author{\blu{Francis Bond}   \\ 
\normalsize  \textbf{Palacký University}\\
\normalsize  \url{https://fcbond.github.io/}\\
\normalsize  \texttt{bond@ieee.org}}
\MyLogo{JMORF (2026)}
\renewcommand{\logo}{#2}
\hypersetup{
   pdfinfo={
     Author={Francis Bond},
     Title={#1: #2},
     Subject={JMORF/MORS: Morpho-Syntax},
     Keywords={Syntax, Morphology, HPSG, Unification, Constructions},
     License={CC BY 4.0}
   }
 }

\date{#1
  \\ Location: KB-1.11}
}
\usepackage{hyperref}
\hypersetup{
     colorlinks,
     linkcolor={blue!50!black},
     citecolor={red!50!black},
     urlcolor={blue!80!black}
   }



\usepackage{xcolor}
\usepackage{graphicx}
\newcommand{\blu}[1]{\textcolor{blue}{#1}}
\newcommand{\grn}[1]{\textcolor{green}{#1}}
\newcommand{\hide}[1]{\textcolor{white}{#1}}
\newcommand{\emp}[1]{\textcolor{red}{#1}}
\newcommand{\txx}[1]{\textbf{\textcolor{blue}{#1}}}
\newcommand{\lex}[1]{\textbf{\mtcitestyle{#1}}}


% \renewcommand{\labelitemi}{\textcolor{violet}{‣}}
% \renewcommand{\labelitemii}{\textcolor{purple}{‣}}

\newcommand{\subhead}[1]{\noindent\textbf{#1}\\[5mm]}

\newcommand{\Bad}{\emp{\raisebox{0.15ex}{\ensuremath{\mathbf{\otimes}}}}}
\newcommand{\bad}[1]{*\eng{#1}}

\newcommand{\com}[1]{\hfill (\emp{#1})}%

\usepackage{relsize,xspace}
\newcommand{\into}{\ensuremath{\rightarrow}\xspace}
\newcommand{\tot}{\ensuremath{\leftrightarrow}\xspace}
\usepackage{url}
\newcommand{\lurl}[1]{\MyLogo{\url{#1}}}

\usepackage{langsci-gb4e}
\let\eachwordone=\slshape
%%% langsci-gb4e defines \thexnumi as \@roman\@xsi{xnumi}, which only works
%%% inside an exe environment; some decks use xlisti on its own.  Restore
%%% the plain gb4e definition so that keeps working.
\makeatletter
\def\thexnumi{\@xsi{xnumi}}
\makeatother
\newcommand{\lx}[1]{\textbf{\mtciteform{#1}}}
\newcommand{\ix}{\ex\slshape}



\newcommand{\ent}{\ensuremath{\Rightarrow}\xspace}
\newcommand{\ngv}{\ensuremath{\not\Rightarrow}\xspace}
%\usepackage{times}
%\usepackage{nttfoilhead}
\newcommand{\myslide}[1]{\foilhead[-25mm]{\raisebox{12mm}[0mm]{\emp{#1}}}\MyLogo{\logo}}
\newcommand{\myslider}[1]{\rotatefoilhead[-25mm]{\raisebox{12mm}[0mm]{\emp{#1}}}}
%\newcommand{\myslider}[1]{\rotatefoilhead{\raisebox{-8mm}{\emp{#1}}}}

\newcommand{\section}[1]{\myslide{}{\begin{center}\Huge \emp{#1}\end{center}}}



\usepackage[lyons,j,e,k]{mtg2e}
%\renewcommand{\mtcitestyle}[1]{\textcolor{-red!75!green!50}{\textsl{#1}}}
\renewcommand{\mtcitestyle}[1]{\textcolor{teal}{\textsl{#1}}}
\newcommand{\cs}{\mtciteform}
\newcommand{\iz}[1]{\texttt{\textup{#1}}}
\newcommand{\gm}{\textsc}
\usepackage[normalem]{ulem}
\newcommand{\ul}{\uline}
\newcommand{\ull}{\uuline}
\newcommand{\wl}{\uwave}
\newcommand{\vs}{\ensuremath{\Leftrightarrow}~}
%%%
%%% Bibliography
%%%
\usepackage{natbib}
%\usepackage{url}
\usepackage{bibentry}



\usepackage[utf8]{inputenc}
%%% trees
\usepackage{multicol}
%%% pst-node: one example in lec-12-control draws a connector between
%%% two points with \rnode/\ncbar.  This used to come in with rtrees.
\usepackage{pst-node}
\usepackage{forest}
\useforestlibrary{linguistics}
\forestapplylibrarydefaults{linguistics}

%%% \makebox's optional arguments are read as AVM delimiters by
%%% langsci-avm, so hide them behind a macro with no visible brackets.
\newcommand{\lbox}[1]{\makebox[0mm][l]{#1}}

%%% Tree styles.
%%%   words  -- trees whose leaves are actual words (italicised, aligned)
%%%   cat    -- trees whose leaves are categories or feature structures
\forestset{
  words/.style={for tree={s sep=7mm, l sep=5mm,
                          if n children=0{font=\itshape, tier=terminal}{}}},
  cat/.style={for tree={s sep=7mm, l sep=5mm}},
}
\newcommand{\wtree}[1]{\begin{forest} words #1 \end{forest}}
\newcommand{\ctree}[1]{\begin{forest} cat #1 \end{forest}}
%%% the rule being applied at this step of a derivation
\newcommand{\step}[1]{\begin{center}\emp{#1}\end{center}}

%%% AVMs (Language Science Press)
\usepackage{langsci-avm}

%%% Macros the decks use, carried over from headx.tex
\newcommand{\blank}{\rule{3em}{1pt}\xspace}
\newcommand{\ft}[1]{\textsc{#1}}
\newcommand{\idn}{\mtciteform}    %%% Indonesian
\newcommand{\prd}[1]{\textbf{#1}}
\newcommand{\typ}[1]{\textit{#1}}
%%% \val: provide before renewing, in case a package defines it.
\providecommand{\val}[1]{}
\renewcommand{\val}[1]{\textit{#1}}
%%% Manning's \avmbox, still used in tree labels and in maths where a
%%% bare langsci \tag would not work.  \avm{\tag{...}} gives the same
%%% boxed tag as the ones inside feature structures.
\newcommand{\avmbox}[1]{\mbox{\avm{\tag{#1}}}}

%%% Ring a value in a dotted oval.  These used pstricks \ovalnode,
%%% which came in with rtrees; tikz (via forest) does the same job.
\newcommand{\OV}[1]{\tikz[baseline]{\node[draw,dotted,red,ellipse,
          inner sep=1pt,outer sep=0pt]{#1};}}
\newcommand{\OVB}[1]{\tikz[baseline]{\node[draw,dotted,blue,ellipse,
          inner sep=1pt,outer sep=0pt]{#1};}}

%%% From CSLI book
\newcommand{\mc}{\multicolumn}
\newcommand{\HD}{\textbf{H}\xspace}
\newcommand{\el}{< >}       %%% empty list; langsci-avm needs the space

\newcommand{\A}{\noindent\textbf{A}: }
\newcommand{\Q}{\noindent\textbf{Q}: }

%%% Question and answer on two pages, written once:
%%%   \qaslide{Title}{... question ... \answer{... answer ...}}
%%% prints the slide twice, first without the \answer parts and then
%%% with them, so the answer appears "on click".  The example counter
%%% is reset so both pages show the same example numbers.
\newif\ifshowanswer
\newcommand{\answer}[1]{\ifshowanswer#1\fi}
\newcounter{qaxnumi}
\newcommand{\qaslide}[2]{%
  \setcounter{qaxnumi}{\value{xnumi}}%
  \showanswerfalse\myslide{#1}#2%
  \setcounter{xnumi}{\value{qaxnumi}}%
  \showanswertrue\myslide{#1}#2%
  \showanswerfalse}
%\newcommand{\C}{\noindent\textbf{C}: }



\begin{document}

\header{Lecture 8}{The Structure of the Lexicon II}{Lexical Rules}
\maketitle



\myslide{Overview}
\begin{enumerate}
\item \txx{Default Inheritance}: how `\,/\,' constraints combine
\item \txx{Lexical Rules}: the Plural Noun LR with everything it
  inherits, licensing \lex{cats}, constant lexemes, \ft{arg-st} and the
  ARP, the feature \ft{form}
  \\ Activity: the word structure for \lex{writers}
\item \txx{Appendix}: the other lexical rules so far, for reference
\end{enumerate}
\bigskip
Then: the passive (a lexical rule at work).

%%%%%%%%%%%%%%%%%%%%%%%%%%%%%%%%%%%%%%%%%%%%%%%%%%%%%%%%%%%%%%%%%%%%%%%%
%%% Converted from Emily Bender's slides (ch08a-uw.pdf, ch08b-uw.pdf),
%%% with permission, and updated to follow her 2026 decks
%%% (1026_lexicaltypes.pdf, 1028_lexicalrules.pdf); the %%% pN comments
%%% give the page in those decks.
%%%%%%%%%%%%%%%%%%%%%%%%%%%%%%%%%%%%%%%%%%%%%%%%%%%%%%%%%%%%%%%%%%%%%%%%
\section{Default Inheritance}

%%% p22
\qaslide{Default Inheritance}{
Why do we have default inheritance?
\answer{%
\begin{itemize}
\item[A] \emp{Generalizations with exceptions are common:}
  \begin{itemize}
  \item Most nouns in English aren't marked for \ft{case}, but
    pronouns are.
  \item Most verbs in English only distinguish two agreement
    categories (\val{3sing} and \val{non-3sing}), but \lex{be} distinguishes
    more.
  \item Most prepositions in English are transitive, but \lex{here} and
    \lex{there} are intransitive.
  \item Most nominal words in English are 3rd person, but some
    (all of them pronouns) are 1st or 2nd person.
  \item Most proper nouns in English are singular, but some
    (mountain range names, sports team names) are plural.
  \end{itemize}
\end{itemize}
\begin{small}
  Note: defaults are not necessary, and currently implemented grammars
  do not use default inheritance, though some work is looking into
  adding it.
\end{small}}
}

%%% p23--28
\newcommand{\dftrow}[3]{%
\raggedright #1 &
\raggedright #2 &
\raggedright #3\tabularnewline}
\qaslide{Default Inheritance, Technicalities}{
\vspace*{1em}
\begin{tabular}{@{}p{0.28\textwidth}p{0.28\textwidth}p{0.32\textwidth}@{}}
\dftrow{If a type says \avm{[ arg-st & / < NP > ]},}%
       {and one of its subtypes says \avm{[ arg-st & < > ]},}%
       {then the \ft{arg-st} value of instances of the subtype is
         \avm{< >}.}
\dftrow{\vspace*{3em}If a type says \avm{[ arg-st & < NP > ]},}%
       {\vspace*{3em}and one of its subtypes says \avm{[ arg-st & < > ]},}%
       {\vspace*{3em}\answer{then this subtype can have no instances,
           since they would have to satisfy contradictory constraints.}}
\end{tabular}
\answer{\par\bigskip
  \emp{In plain English: without the `\,/\,' the constraint is
    inviolable.  Nothing can be both \avm{< NP >} and \avm{< >}, so
    the subtype is empty.}}
}

%%% p29
\qaslide{Default Inheritance, More Technicalities}{
\begin{itemize}
\item If a type says \avm{[ mod & / < S > ]}, and one of its subtypes says
  \avm{[ mod & < [ spr & < NP > ] > ]}, then the \ft{mod} value of
  instances of the subtype is what?
\end{itemize}
\answer{%
\begin{center}
  \avm{[ mod & < [ head & / \type{verb} \\
                   spr & < NP > \\
                   comps & / < > ] > ]}
\end{center}
\begin{itemize}
\item That is, default constraints are `pushed down'
\item \emp{S abbreviates
    \avm{[ head & \type{verb} \\ spr & < > \\ comps & < > ]}, so the
    subtype only overrides the part it mentions (\ft{spr}).  The rest
    of the default survives, still as a default.}
\end{itemize}}
}

%%% p30
\qaslide{Question on Default Inheritance}{
\begin{itemize}
\item[Q] Can a grammar rule override a default
  constraint on a word?
\answer{\item[A] \emp{No. Defaults are all `cached out' in the
  lexicon.}
\item Words as used to build sentences have only
  inviolable constraints.
  \begin{itemize}
  \item i.e.\ nodes of type \val{word} in your tree should never be
    shown with a defeasible constraint (no `\,/\,' in a tree)
  \end{itemize}}
\end{itemize}
}

\section{Lexical Rules}

%%% Grey box over a part of an AVM that is not yet revealed (her
%%% build-up slides cover parts of the AVM with grey rectangles).
\providecommand{\gry}[1]{{\setlength{\fboxsep}{0pt}\colorbox{lightgray}{\phantom{#1}}}}
%%% Boxes holding hidden AVM parts that contain & (these cannot be
%%% passed through \gry inside an \avm directly).
\ifdefined\gryA\else\newsavebox{\gryA}\fi
\ifdefined\gryB\else\newsavebox{\gryB}\fi
\ifdefined\gryC\else\newsavebox{\gryC}\fi


%%% p4
\myslide{Overview}
\begin{itemize}
\item Recap: the three subtypes of \val{l-rule}
\item The Plural Noun LR, and what it inherits
\item Example: licensing \lex{cats}
\item Advice on writing lexical rules
\item Constant lexemes
\item \ft{arg-st} \& ARP; the feature \ft{form}
\item Activity: the word structure for \lex{writers}
\end{itemize}

%%% p14
\myslide{Recap: Three Subtypes of \textit{l-rule}}
\begin{center}
\begin{small}
  \begin{forest}
    cat, for tree={font=\itshape, l sep=3mm}
    [{l-rule} [{i-rule}] [{d-rule}] [{pi-rule}]]
  \end{forest}

  \medskip
  \val{l-rule} :
  \avm{[ input & \type{l-sequence} < X , [ sem & / \2 ] > \\
         output & \type{l-sequence} < Y , [ sem & / \2 ] > ]}

  \medskip
  \val{i-rule} :
  \avm{[ input & < X , [\type*{lexeme}
                         syn & \3 \\
                         arg-st & \tag{A} ] > \\
         output & < Y , [\type*{word}
                          syn & \3 \\
                          arg-st & \tag{A} ] > ]}
  \qquad
  \val{d-rule} :
  \avm{[ input & < X , [\type*{lexeme}
                         syn & / \3 ] > \\
         output & < Y , [\type*{lexeme}
                          syn & / \3 ] > ]}
\end{small}
\end{center}
\begin{small}
\emp{\val{l-rule}: the output keeps the input's meaning, by default.
  \val{i-rule}: the output is a \val{word} with the same \ft{syn} and
  \ft{arg-st}, no exceptions.
  \val{d-rule}: the output is another \val{lexeme} with the same
  \ft{syn}, by default.}
\end{small}

%%% p15
\myslide{Plural Noun LR}
\begin{center}
  \avm{[\type*{i-rule}
    input & < \1 , \type{cntn-lxm} > \\
    output & < F$_{NPL}$(\1) , [\type*{word}
                syn & [ head & [ agr & [ num & pl ] ] ] ] > ]}
\end{center}
\emp{Take any count noun lexeme \avm{\1}, add the plural ending to
  its form, and make it a word whose number is plural.  That is all
  the rule itself says; everything else comes from the types.}

%%% p25
\myslide{Practicalities --- Applying Lexical Rules}
\begin{itemize}
\item \ft{input} is a family of lexical sequences.
\item \ft{output} is another family of lexical sequences.
  \begin{itemize}
  \item \ldots usually a smaller family
  \item \ldots usually a disjoint one
  \end{itemize}
\item The only differences between the families are those
  stipulated in the rule (or the rule's type).
\item Similarities are handled by the constraints on \val{l-rule}
  and its subtypes.
\item If we've written the LRs correctly, nothing is left
  underconstrained.
\end{itemize}

%%% p26
\myslide{Example: Lexical Entry for \textit{cat}}
\begin{center}
  \avm{< \textnormal{cat} , [\type*{cntn-lxm}
    sem & [ index & $k$ \\
            restr & < [ reln & cat \\ inst & $k$ ] > ] ] >}
\end{center}
\emp{The entry gives only the form, the type and the meaning.}

%%% p27--30
\myslide{Example: \textit{cat}, with Inheritance}
\begin{center}
\sbox{\gryA}{\avm{[\type*{noun} agr & \7 [ per & 3rd ] ]}}
\sbox{\gryB}{\avm{< DP[ count & + \\ agr & \7 ] >}}
  \avm{< \textnormal{cat} , [\type*{cntn-lxm}
    syn & [ \gry{head} & \gry{\usebox{\gryA}} \\
            val & [ \gry{spr} & \gry{\usebox{\gryB}} ] ] \\
    sem & [ \gry{mode} & \gry{ref} \\
            index & $k$ \\
            restr & < [ reln & cat \\ inst & $k$ ] > ] \\
    \gry{arg-st} & \gry{\avm{< X >}} ] >}
\end{center}
\emp{The types fill in the rest.  Which type adds what?}

\myslide{Example: \textit{cat}, with Inheritance --- \val{cntn-lxm}}
\begin{center}
\sbox{\gryA}{\avm{[\type*{noun} agr & \7 [ per & 3rd ] ]}}
  \avm{< \textnormal{cat} , [\type*{cntn-lxm}
    syn & [ \gry{head} & \gry{\usebox{\gryA}} \\
            val & [ spr & < \gry{DP}[ count & + \\ \gry{agr} & \gry{\avm{\7}} ] > ] ] \\
    sem & [ \gry{mode} & \gry{ref} \\
            index & $k$ \\
            restr & < [ reln & cat \\ inst & $k$ ] > ] \\
    \gry{arg-st} & \gry{\avm{< X >}} ] >}
\end{center}
\emp{\val{cntn-lxm}: the specifier is [\ft{count} +].}

\myslide{Example: \textit{cat}, with Inheritance --- \val{cn-lxm}}
\begin{center}
  \avm{< \textnormal{cat} , [\type*{cntn-lxm}
    syn & [ head & [\type*{noun} agr & \gry{\avm{\7}} [ per & 3rd ] ] \\
            val & [ spr & < DP[ count & + \\ \gry{agr} & \gry{\avm{\7}} ] > ] ] \\
    sem & [ mode & ref \\
            index & $k$ \\
            restr & < [ reln & cat \\ inst & $k$ ] > ] \\
    arg-st & < X > ] >}
\end{center}
\emp{\val{cn-lxm}: \val{noun}, [\ft{per} 3rd], a DP specifier,
  [\ft{mode} ref] and one argument.  \txx{Defeasible constraints are
    inviolable in lexical entries!}  The type says [\ft{mode} / ref];
  nothing in the entry overrides it, so for \lex{cat} it is simply
  [\ft{mode} ref].}

\myslide{Example: \textit{cat}, with Inheritance --- \val{infl-lxm}}
\begin{center}
  \avm{< \textnormal{cat} , [\type*{cntn-lxm}
    syn & [ head & [\type*{noun} agr & \7 [ per & 3rd ] ] \\
            val & [ spr & < DP[ count & + \\ agr & \7 ] > ] ] \\
    sem & [ mode & ref \\
            index & $k$ \\
            restr & < [ reln & cat \\ inst & $k$ ] > ] \\
    arg-st & < X > ] >}
\end{center}
\emp{\val{infl-lxm}: the SHAC identifies the noun's \ft{agr} with
  its specifier's, \avm{\7}.}

%%% p31
\myslide{Plural Noun LR}
\begin{center}
  \avm{[\type*{i-rule}
    input & < \1 , \type{cntn-lxm} > \\
    output & < F$_{NPL}$(\1) , [\type*{word}
                syn & [ head & [ agr & [ num & pl ] ] ] ] > ]}
\end{center}
\emp{Now feed \lex{cat} to the rule.}

%%% p32--34
\myslide{Licensing \textit{cats}: the Plural Noun LR's Own Contributions}
\begin{center}
\scalebox{0.7}{%
  \avm{[\type*{i-rule}
    input & < \1\textnormal{cat} , [\type*{cntn-lxm}
       syn & \3 [ head & [\type*{noun} agr & \7 [ per & 3rd ] ] \\
                  val & [ spr & < DP[ count & + \\ agr & \7 ] > ] ] \\
       sem & \2 [ mode & ref \\
                  index & $k$ \\
                  restr & < [ reln & cat \\ inst & $k$ ] > ] \\
       arg-st & \emp{\avm{\tag{B}}} < X > $\oplus$ \tag{C} < > ] > \\
    output & < \emp{F$_{NPL}$(\avm{\1})} , [\type*{\emp{word}}
       syn & \3 [ head & [ agr & \emp{[num & pl]} ] \\
                  val & [ spr & \emp{\avm{\tag{B}}} \\ comps & \tag{C} ] ] \\
       sem & \2 \\
       arg-st & \tag{B} $\oplus$ \tag{C} ] > ]}}
\end{center}
\vspace{-1ex}
\emp{In blue: what the rule and its output type add; the rest is copied or inherited.}

%%% p35
\myslide{\textit{cats}: The (Family of) Lexical Sequence(s)}
\begin{center}
  \avm{< \textnormal{cats} , [\type*{word}
    syn & [ head & [\type*{noun} agr & \7 \type{3pl} ] \\
            val & [ spr & \tag{B} < DP[ count & + \\ agr & \7 ] > \\
                    comps & < > ] ] \\
    sem & [ mode & ref \\
            index & $k$ \\
            restr & < [ reln & cat \\ inst & $k$ ] > ] \\
    arg-st & \tag{B} ] >}
\end{center}
\emp{Strictly a family: the DP in \ft{spr} is not fully resolved, so
  this stands for every lexical sequence that fills it in.}

%%% p36
\myslide{Practicalities --- Writing Lexical Rules}
\begin{itemize}
\item Determine the type of the LR.
\item Determine the class of possible inputs.
\item Determine what should change.
  \begin{small}
  \begin{itemize}
  \item If \ft{input} and \ft{output} values are identified (by default or otherwise) and
    only \ft{output} value is mentioned, then\ldots
    \\ \emp{information is added.}  e.g.\ Plural Noun LR
    \\[1ex] (Lexical sequences incompatible with that value are not possible inputs)
  \item If \ft{input} and \ft{output} values are identified by default, but different values
    are given on the \ft{input} and \ft{output} of the rule, then\ldots
    \\ \emp{information is changed.}  e.g.\ Agent Nominalization LR
  \item If \ft{input} and \ft{output} values are identified by an inviolable constraint, but
    different values are given on the \ft{input} and \ft{output} of the rule, then\ldots
    \\ \emp{there is no well-formed input: the rule can never apply.}
  \end{itemize}
  \end{small}
\end{itemize}

%%% p37
\qaslide{Constant Lexemes}{
\begin{itemize}
\item What kinds of words are constant lexemes
  in our grammar?
  \answer{\\ \emp{Adjectives, adverbs, prepositions, determiners,
      conjunctions: words that do not inflect in English.}}
\item Why do we need a rule for these words?
  \answer{\\ \emp{They need to go `from' \val{lexeme} to \val{word},
      and only words appear in trees.}}
\item What would be an alternative analysis?
  \answer{\\ \emp{Put the sub-hierarchy for constant lexemes under
      \val{word} instead of under \val{lexeme}.}}
\end{itemize}
}

%%% p38--39
\qaslide{Constant Lexeme LR}{
\begin{center}
  \avm{[\type*{i-rule}
    input & < \1 , \type{const-lxm} > \\
    output & [ first & \1 ] ]}
  \qquad i.e.\qquad
  \avm{[\type*{i-rule}
    input & < \1 , \type{const-lxm} > \\
    output & < \1 , [ ] > ]}
\end{center}
\emp{The form does not change, and the rule says nothing else.  The
  type \val{i-rule} does the work: it copies \ft{syn} and \ft{arg-st}
  and makes the output a \val{word}.}
\begin{itemize}
\item What keeps this from applying to, say, verb
  lexemes?
  \answer{\\ \emp{The input must be a \val{const-lxm}.  Verb lexemes
      are \val{infl-lxm}s, a different branch of the hierarchy.}}
\item Why is this an \val{i-rule}?
  \answer{\\ \emp{It takes a \val{lexeme} to a \val{word}, which is
      what inflectional rules do, even though nothing is inflected.}}
\end{itemize}
}

%%% p49
\qaslide{\ft{arg-st} \& ARP}{
\begin{itemize}
\item Given the ARP, what do we need to specify
  about the valence properties of words?
  \answer{\\ \emp{Only the length of the \ft{spr} list, and the SHAC
      already does that: nouns and verbs have exactly one specifier.
      The ARP then fills \ft{spr} and \ft{comps} from \ft{arg-st}.}}
\item Why isn't the ARP a constraint on the type
  \val{lexeme}?
  \answer{\\ \emp{Some derivational rules change the argument
      structure (e.g.\ Agent Nominalization).  If \ft{spr} and
      \ft{comps} were already linked to \ft{arg-st} in the lexeme,
      those rules would have to change all three.}}
\end{itemize}
}

%%% p50
\qaslide{The Feature \ft{form}}{
\begin{itemize}
\item Different inflected forms of verbs show
  up in different syntactic environments.
  \\ Examples?
  \answer{\\ \emp{\eng{I eat}; \eng{Eat!}; \eng{I was eating};
      \eng{I have eaten}}}
\item These different forms are syntactically
  distinguished by the feature \ft{form}, as
  assigned by lexical rules.
\item \ft{form} is also useful in our analyses of
  coordination and PP selection.
\end{itemize}
}

%%% p51--52
\qaslide{What Rules These Out?}{
\begin{exe}
\ix \bad{Kim eat pizza.}
  \answer{\\ \emp{\eng{eat} is either Non-3rd-Singular (does not
      agree with \eng{Kim}) or Base Form (not finite, so it cannot
      head a sentence)}}
\ix \bad{Kim seems to eats pizza.}
  \answer{\\ \emp{Spoiler: \eng{to} will be an auxiliary verb that
      takes a [\ft{form} base] complement (Chapter 12), so \eng{eats}
      does not fit}}
\ix \bad{Dana helped Leslie pack and moved.}
  \answer{\\ \emp{\eng{helped} selects a [\ft{form} base] complement;
      the Coordination Rule makes the conjuncts and the mother share
      \ft{form}, so \eng{moved} cannot join \eng{pack}}}
\ix \bad{Kim relies for Sandy.}
  \answer{\\ \emp{\eng{relies} requires a PP complement headed by a
      preposition with [\ft{form} on]}}
\ix \bad{Dana walked and Kim.}
  \answer{\\ \emp{Used to be fine, since the \ft{val} values match;
      now the Coordination Rule requires the conjuncts to share
      \ft{form}, and \eng{Kim} has none}}
\end{exe}
}

%%% p53
\myslide{Review: Lexical Terminology}
\begin{small}
\begin{itemize}
\item \txx{Lexeme}: a type in our grammar; its subtypes state
  generalizations about combinations of constraints
\item \txx{Lexical entry}: what we write in the lexicon; it licenses
  a family of lexical sequences
\item \txx{Lexical sequence}: a form paired with a fully resolved
  feature structure of type \val{lexeme} or \val{word}
\item \txx{Family of lexical sequences}: a form paired with an
  underspecified feature structure, resolvable in many ways
\item \txx{Lexical rule}: what we write to derive new lexical
  sequences from old ones
\item \txx{Lexical rule instantiation}: a fully resolved \val{l-rule}:
  one rule, one input, one output
\item \txx{Word structure}: the bottom of the tree: a form with a
  fully resolved \val{word} feature structure over it
\end{itemize}
\end{small}

%%% p54--56: in-class activity
\qaslide{Activity: Arriving at a Word Structure}{
Write the word structure for \lex{writers} as it appears in the tree
for the sentence
\begin{exe}
\ix Those writers need inspiration.
\end{exe}
starting from the lexical entry for \lex{write}:
\begin{center}
  \avm{< \textnormal{write} , [\type*{stv-lxm}
    sem & [ index & $s$ \\
            restr & < [ reln & write \\ sit & $s$ \\
                        writer & $i$ \\ written & $j$ ] > ] \\
    arg-st & < X$_i$, Y$_j$ > ] >}
\end{center}
\answer{\emp{Which lexical rules do we need, and in which order?
    Agent Nominalization LR (\lex{write} $\to$ \lex{writer}, a
    \val{d-rule}), then Plural Noun LR (\lex{writer} $\to$
    \lex{writers}, an \val{i-rule}).}}
}

\myslide{Step 1: Agent Nominalization LR}
\begin{center}
\scalebox{0.8}{%
  \avm{< \textnormal{writer} , [\type*{cntn-lxm}
    syn & [ head & [\type*{noun} agr & \1 [ per & 3rd ] ] \\
            val & [ spr & < DP[ count & + \\ agr & \1 ] > ] ] \\
    sem & [ mode & ref \\
            index & \emp{$i$} \\
            restr & < [ reln & write \\ sit & $s$ \\
                        writer & $i$ \\ written & $j$ ] > ] \\
    arg-st & < DP ( , PP$_j$[ form & of ] ) > ] >}}
\end{center}
\emp{The rule adds \textit{-er}, changes the type to \val{cntn-lxm}
  and makes the index the writer, $i$; the \ft{restr} is unchanged.
  The noun types add the rest.  The verb's object $j$ survives only
  as an optional \eng{of}-PP.}

\myslide{Step 2: Plural Noun LR}
\begin{center}
\scalebox{0.8}{%
  \avm{< \textnormal{writers} , [\type*{word}
    syn & [ head & [\type*{noun} agr & \1 \type{3pl} ] \\
            val & [ spr & \tag{B} < DP[ count & + \\ agr & \1 ] > \\
                    comps & \tag{C} < > ] ] \\
    sem & [ mode & ref \\
            index & $i$ \\
            restr & < [ reln & write \\ sit & $s$ \\
                        writer & $i$ \\ written & $j$ ] > ] \\
    arg-st & \tag{B} $\oplus$ \tag{C} ] >}}
\end{center}
\emp{Add \textit{-s} and make it a \val{word} with [\ft{num} pl].
  Now the ARP applies: \ft{spr} is the DP and \ft{comps} is empty,
  since here \lex{writers} has no \eng{of}-PP.}

\myslide{Activity: The Word Structure for \lex{writers}}
\begin{center}
\scalebox{0.7}{%
\ctree{[{\avm{[\type*{word}
    syn & [ head & [\type*{noun} agr & \1 \type{3pl} ] \\
            val & [ spr & < \2 DP[ count & + \\ agr & \1 ] > \\
                    comps & < > ] ] \\
    sem & [ mode & ref \\
            index & $i$ \\
            restr & < [ reln & write \\ sit & $s$ \\
                        writer & $i$ \\ written & $j$ ] > ] \\
    arg-st & < \2 > ]}}
  [{\eng{writers}}]]}}
\end{center}
\emp{The form \eng{writers} under a feature structure of type
  \val{word}, licensed by two lexical rules from the entry for
  \lex{write}.}


\myslide{Summary}
\begin{itemize}
\item A lexical rule says only what changes; its type (\val{i-rule},
  \val{d-rule}, \val{pi-rule}) and the input's lexeme type supply the
  rest, and defaults let a rule override what it needs to
\item Licensing a word: lexical entry $\rightarrow$ lexical rule(s)
  $\rightarrow$ word structure at the bottom of the tree
\item Constant lexemes still pass through an \val{i-rule} to become
  words
\item The ARP links \ft{arg-st} to \ft{spr} and \ft{comps} on words;
  \ft{form} records which inflected form a word is
\item Next: the passive, a derivational rule that rearranges
  \ft{arg-st}
\end{itemize}

\section{Appendix: The Lexical Rules So Far}

%%% p40--48: the other lexical rules, from SWB Chapter 9
\myslide{Other Lexical Rules (So Far)}
\begin{itemize}
\item Inflectional (\val{i-rule}): Singular Noun, 3rd-Singular Verb,
  Non-3rd-Singular Verb, Past-Tense Verb, Base Form
\item Derivational (\val{d-rule}): Present Participle, Past
  Participle, Agent Nominalization
\item Each one is in the grammar summary in Chapter 9.  For reference, on the next
  slides, read each rule and ask: what is the input, what changes, and
  where does the rest come from?
\end{itemize}

\myslide{Singular Noun Lexical Rule}
\begin{center}
  \avm{[\type*{i-rule}
    input & < \1 , \type{cn-lxm} > \\
    output & < \1 , [ syn & [ head & [ agr & [ num & sg ] ] ] ] > ]}
\end{center}
\emp{Any common noun lexeme, count or mass, gives a singular word
  with the same form.  Compare the Plural Noun LR, whose input is
  only \val{cntn-lxm}: mass nouns have no plural.}

\myslide{3rd-Singular Verb Lexical Rule}
\begin{center}
  \avm{[\type*{i-rule}
    input & < \3 , [\type*{verb-lxm} sem & [ restr & \tag{A} ] ] > \\
    output & < F$_{3SG}$(\3) ,
              [ syn & [ head & [ form & fin \\ agr & \type{3sing} ] ] \\
                sem & [ restr & \tag{A} $\oplus$ \ldots ] \\
                arg-st & < [ case & nom ], \ldots > ] > ]}
\end{center}
\emp{Add \textit{-s}; the word is finite and 3rd singular, its
  subject (the first argument) is nominative, and the \ft{restr}
  gets something extra (the tense).}

\myslide{Non-3rd-Singular Verb Lexical Rule}
\begin{center}
  \avm{[\type*{i-rule}
    input & < \1 , [\type*{verb-lxm} sem & [ restr & \tag{A} ] ] > \\
    output & < \1 ,
              [ syn & [ head & [ form & fin \\ agr & \type{non-3sing} ] ] \\
                sem & [ restr & \tag{A} $\oplus$ \ldots ] \\
                arg-st & < [ case & nom ], \ldots > ] > ]}
\end{center}
\emp{The same, but the form stays as it is and the agreement is
  anything but 3rd singular: \eng{I/you/we/they walk}.}

\myslide{Past-Tense Verb Lexical Rule}
\begin{center}
  \avm{[\type*{i-rule}
    input & < \3 , [\type*{verb-lxm} sem & [ restr & \tag{A} ] ] > \\
    output & < F$_{PAST}$(\3) ,
              [ syn & [ head & [ form & fin ] ] \\
                sem & [ restr & \tag{A} $\oplus$ \ldots ] \\
                arg-st & < [ case & nom ], \ldots > ] > ]}
\end{center}
\emp{Add \textit{-ed}; finite, nominative subject, but no agreement
  constraint at all: \eng{I/you/she/they walked}.}

\myslide{Base Form Lexical Rule}
\begin{center}
  \avm{[\type*{i-rule}
    input & < \1 , \type{verb-lxm} > \\
    output & < \1 , [ syn & [ head & [ form & base ] ] ] > ]}
\end{center}
\emp{Same form, [\ft{form} base]: not finite, so no nominative
  subject and no agreement.  This is the verb in \eng{Eat!}, \eng{to
    eat} and \eng{will eat}.}

\myslide{Present Participle Lexical Rule}
\begin{center}
  \avm{[\type*{d-rule}
    input & < \3 , [\type*{verb-lxm} sem & [ restr & \tag{A} ] \\
                                    arg-st & \tag{B} ] > \\
    output & < F$_{PRP}$(\3) ,
              [\type*{part-lxm}
                syn & [ head & [ form & prp ] ] \\
                sem & [ restr & \tag{A} $\oplus$ \ldots ] \\
                arg-st & \tag{B} ] > ]}
\end{center}
\emp{A \val{d-rule}: \lex{eat} gives the lexeme \lex{eating}, of
  type \val{part-lxm}, with the same arguments.  \val{part-lxm} is a
  \val{const-lxm}, so the Constant Lexeme LR then makes it a word.}

\myslide{Past Participle Lexical Rule}
\begin{center}
  \avm{[\type*{d-rule}
    input & < \3 , [\type*{verb-lxm} sem & [ restr & \tag{A} ] \\
                                    arg-st & \tag{B} ] > \\
    output & < F$_{PSP}$(\3) ,
              [\type*{part-lxm}
                syn & [ head & [ form & psp ] ] \\
                sem & [ restr & \tag{A} $\oplus$ \ldots ] \\
                arg-st & \tag{B} ] > ]}
\end{center}
\emp{The same, with [\ft{form} psp]: \lex{eat} gives \lex{eaten}, as
  in \eng{I have eaten}.}

\myslide{Agent Nominalization Lexical Rule}
\begin{center}
  \avm{[\type*{d-rule}
    input & < \2 , [\type*{verb-lxm} sem & [ index & $s$ ] \\
                                    arg-st & < X$_i$ , NP$_j$ > ] > \\
    output & < F$_{\textit{-er}}$(\2) ,
              [\type*{cntn-lxm}
                sem & [ index & $i$ ] \\
                arg-st & < Y ( , PP$_j$[ form & of ] ) > ] > ]}
\end{center}
\emp{Add \textit{-er} to a strictly transitive verb lexeme.  The
  result is a count noun lexeme whose index is the verb's subject
  index $i$: a \lex{driver} is the one who drives.  The verb's object
  $j$ can only appear as an \eng{of}-PP: \eng{a builder of bridges}.}

\myslide{Acknowledgments and References}

\begin{itemize}
\item Course design and slides borrow heavily from Emily Bender's course:
\textit{Linguistics 566: Introduction to Syntax for Computational Linguistics}
\\ \url{http://courses.washington.edu/ling566}
\item The slides on lexical types and lexical rules are adapted from
  Emily Bender's, with her permission; the feature structures follow
  \citet{sag:wasow:bender:2003}, \copyright{} 2003 CSLI Publications
\item This lecture complements \citet[Chapter 8]{sag:wasow:bender:2003}
\end{itemize}


%\input{lec-05-questions}

\bibliographystyle{aclnat}
\bibliography{abb,mtg,nlp,ling}
\end{document}

%%% Local Variables: 
%%% coding: utf-8
%%% mode: latex
%%% TeX-PDF-mode: t
%%% TeX-engine: xetex
%%% End: 
